% SHARD-ID: CA-35-QUBIT-1D-BOOST-CURRENT-OBSTRUCTIONS
% SHARD-TITLE: One-Dimensional Qubit Boost-Current Obstructions
% SHARD-SUMMARY: Derives the adjacent-current polynomial for a qubit two-site Pauli density.
% SHARD-SUMMARY: Checks that symmetric fully on-site Hamiltonians and commuting classical densities have zero CA-12 bulk momentum candidate.
% SHARD-SUMMARY: Records why this rules out only the first-moment route to nontrivial Poincare symmetry, not every possible construction.
% SHARD-KEYWORDS: qubit chain, Pauli basis, boost current, adjacent commutator, onsite obstruction, Poincare diagnostic

\section{One-Dimensional Qubit Boost-Current Obstructions}
\label{sec:qubit-1d-boost-current-obstructions}

\subsection{Status, Sources, and Dependencies}

\textbf{Status: checked local derivation.}  This shard turns the CA-12
nearest-neighbour identity into explicit polynomial conditions on the
two-qubit Pauli coefficient matrix.  The obstruction is checked in Julia.

Dependencies: CA-12, CA-34, CONVENTIONS.md (h), (k), (m).

% Source: report/sections/12_lattice_boost_current_1d.tex:24--104 -- local
%   derivation of i[H,K] as the adjacent-density current sum.
% Convention: CONVENTIONS.md (m) -- qubit Pauli coefficient normalization and
%   symmetric on-site split.
% Checked: src/QubitPauliLattice.jl and test/runtests.jl, testset
%   "qubit Pauli nearest-neighbour current obstructions".

\subsection{Adjacent Current in Pauli Coefficients}

Let
\[
  h_j=\sum_{\alpha,\beta=0}^3
    h_{\alpha\beta}\,\sigma_j^\alpha\sigma_{j+1}^\beta.
\]
The only noncommuting overlap between \(h_j\) and \(h_{j+1}\) is at site
\(j+1\).  With
\[
  [\sigma^b,\sigma^c]=2i\sum_{e=1}^3\epsilon_{bce}\sigma^e
  \qquad (b,c\in\{1,2,3\}),
\]
one obtains
\[
\begin{aligned}
  i[h_j,h_{j+1}]
  &=
  -2\sum_{\alpha,\delta=0}^3\sum_{e=1}^3
    \left(\sum_{b,c=1}^3
      h_{\alpha b}h_{c\delta}\epsilon_{bce}
    \right)
    \sigma_j^\alpha\sigma_{j+1}^e\sigma_{j+2}^\delta .
\end{aligned}
\]
Equivalently, define right-leg and left-leg vectors
\[
  R_\alpha=(h_{\alpha1},h_{\alpha2},h_{\alpha3}),\qquad
  L_\delta=(h_{1\delta},h_{2\delta},h_{3\delta}).
\]
Then the current coefficients are the cross products
\[
  i[h_j,h_{j+1}]
  =
  -2\sum_{\alpha,\delta=0}^3
  \sum_{e=1}^3 (R_\alpha\times L_\delta)_e
  \sigma_j^\alpha\sigma_{j+1}^e\sigma_{j+2}^\delta .
\]
Thus the simple current-collapse condition is
\[
  R_\alpha\times L_\delta=0
  \qquad\text{for every }\alpha,\delta.
\]
When this holds, every adjacent-current contribution to the CA-12 bulk momentum
candidate vanishes.

\subsection{Fully On-Site Terms}

For a genuine translation-invariant on-site density
\[
  A=\sum_{\alpha=0}^3 a_\alpha\sigma^\alpha,
\]
CONVENTIONS.md (m) embeds it into the two-site bond tier by the symmetric split
\[
  h_{\mathrm{onsite}}
  =
  \frac12(A\otimes I+I\otimes A).
\]
Then the right-leg and left-leg vectors are all parallel in the only places
where they overlap, so
\[
  i[(h_{\mathrm{onsite}})_{12},(h_{\mathrm{onsite}})_{23}]=0.
\]
The Julia test checks this directly for a generic real linear combination of
\(I,X,Y,Z\).

Consequently the first-moment boost ansatz gives
\[
  i[H,K]=0
\]
in the translation-invariant bulk.  If the target is a nontrivial Poincare
representation with \(i[H,K]=P\) and \(i[P,K]=H\), this first-moment route
cannot realise it unless the energy itself is trivial after the relevant
compression or subtraction.

\subsection{Other Immediate Rejections}

The same current-collapse test catches commuting nearest-neighbour density
families.  A pure diagonal classical interaction such as \(ZZ\) has
\[
  [h_j,h_{j+1}]=0
\]
because the overlapping Pauli factor is the same diagonal matrix.  Such a model
may still be a well-defined lattice Hamiltonian, but it fails the first
nontrivial bulk-symmetry diagnostic: the boost-current candidate supplies no
spatial momentum.

By contrast, the tested transverse-Ising-style density with a \(ZZ\) bond and a
symmetric \(X\)-field split has nonzero adjacent current.  This does not prove
Poincare invariance; it only passes the first obstruction.

\subsection{The Next Residual}

The next necessary Poincare relation is translation conservation:
\[
  i[H,P]=0.
\]
For the local candidate \(P=\sum_j C_j\), \(C_j=i[h_j,h_{j+1}]\), the first
bulk density to inspect is
\[
  A_j=i[h_j+h_{j+1},C_j].
\]
If \(A_j\) contains interior Pauli strings that are not a named discrete
divergence or boundary term, the model fails the next finite algebraic filter.

This residual is useful for detecting density-split artifacts.  An asymmetric
encoding of an on-site Hamiltonian,
\[
  h=U\otimes I+I\otimes V,
\]
can create the fake current
\[
  i[h_{12},h_{23}]=I\otimes i[V,U]\otimes I.
\]
But the checked sentinel with \(U=X\), \(V=Z\) has nonzero
\[
  i[h_{12}+h_{23},i[h_{12},h_{23}]],
\]
so it fails translation conservation even though the adjacent current is
nonzero.  A future residual suite should add the next boost relation
\(i[P,K]=\lambda H\) after fixing the allowed density-divergence and boundary
terms.

\subsection{Boundary}

The obstruction is conditional on the density convention.  An asymmetric
encoding of an on-site Hamiltonian as \(A\otimes I+I\otimes B\) can create
spurious adjacent commutators when \(A\) and \(B\) fail to commute.  Such
currents are density-split artifacts unless a separate convention justifies the
split, and the conservation residual above is the first backstop against them.
Passing the obstruction is also far weaker than CA-15 acceptance: boundary
residuals, scaling maps, low-energy sectors, and observable covariance remain
unproved.
